% Abstract (150-200 words)
% TODO: Fill in after completing paper

\textbf{Background:} Small language models (LLMs) benefit from emotion-aware training, but current approaches rely on deep learning embeddings.

\textbf{Method:} We propose using interpretable acoustic features based on speech science principles (prosody, spectral analysis, formants) for emotion-augmented knowledge distillation.

\textbf{Key Finding:} Our 46-dimensional acoustic features achieve \textbf{40.45\% validation loss reduction}, outperforming 256-dimensional WavLM embeddings (27.82\% reduction) by \textbf{17.5\%} with only 18\% of the parameters.

\textbf{Contributions:}
(1) First systematic comparison of speech science vs. deep learning features for emotion-augmented LLM training.
(2) Demonstration that interpretable features outperform black-box representations.
(3) 8.06$\times$ dimensional efficiency with superior performance.

\textbf{Implications:} Challenges the conventional wisdom that deep learning always outperforms hand-crafted features, validating speech science principles for production deployment.
